\documentclass[../../../include/open-logic-section]{subfiles}

\begin{document}

\olfileid[fr]{sth}{choice}{banach}
\olsection{Le paradoxe de Banach--Tarski}

Nous pourrions aussi tenter de justifier
l'axiome du choix, comme Boolos a tenté
de justifier le remplacement, en invoquant
des considérations \emph{extrinsèques}
(voir la \olref[replacement][extrinsic]{sec}).
Après tout, l'adoption du choix a de
nombreuses conséquences souhaitables :
on peut comparer tous les cardinaux,
bien ordonner tous les ensembles,
traiter les cardinaux comme une sorte
particulière d'ordinaux, et ainsi de suite.

On affirme parfois, cependant, que le choix
a des conséquences \emph{indésirables}.
Cette objection s'appuie surtout sur un
résultat de \cite{BanachTarski1924}.

\begin{thm}[Paradoxe de Banach--Tarski
(dans $\ZFC$)]
Toute boule de l'espace euclidien de
dimension trois peut être décomposée en
un nombre fini de parties que l'on peut
réassembler, par rotations et translations,
pour former deux copies de cette boule.\footnote{
Note de l'édition française : l'original
dit seulement «~toute boule~» ;
nous explicitons l'espace euclidien
à trois dimensions de cet énoncé.}
\end{thm}
\noindent
À première vue, c'est étonnant.
Les deux boules ont évidemment
\emph{deux fois} le volume de la boule
initiale. Pourtant, les déplacements
rigides — rotations et translations —
ne changent pas le volume. Le paradoxe
de Banach--Tarski semble donc nous
permettre de faire apparaître de la
matière comme par magie.

Il y a pire.\footnote{Voir
\citet[Theorem 3.12]{Wagon2016}.}
Un raisonnement analogue montre
qu'on peut découper un petit pois
en un nombre fini de parties et
les réassembler, par rotations et
translations, en un objet ayant
la forme et la taille de Big Ben.

Rien de tout cela, cependant, ne
se démontre dans $\ZF$ seule.\footnote{
Le paradoxe de Banach--Tarski peut
néanmoins être démontré à partir
de principes strictement plus faibles
que l'axiome du choix ; voir
\citet[303]{Wagon2016}.}
Il nous faut donc choisir : rejeter
l'axiome du choix, ou apprendre à vivre
avec ce «~paradoxe~».

Nous allons suggérer qu'il faut apprendre
à vivre avec lui. Nous pensons même
qu'il n'a pas grand-chose de paradoxal.
En particulier, nous ne voyons pas
pourquoi il serait plus ou moins
paradoxal que les résultats
suivants :\footnote{
\citet[276--7]{Potter2004},
\citet[16]{Weston2003} et
\citet[31, 308--9]{Wagon2016}
avancent des remarques semblables
à partir d'autres exemples.}
\begin{enumerate}
	\item L'intervalle $(0,1)$ a autant
	de points que $\Real$.
	\\\emph{Démonstration} : considérons
	$\tan(\pi(r-\nicefrac{1}{2}))$.\footnote{
	Note de l'édition française : une
	parenthèse fermante supplémentaire
	figure dans la formule de l'original.}
	\item Une droite et un carré ont
	autant de points.
	\\Voir la \olref[his][set][pathology]{sec}
	et la \olref[his][set][cantorplane]{sec}.
	\item Il existe des courbes qui
	remplissent l'espace.
	\\Voir la \olref[his][set][pathology]{sec}
	et la \olref[his][set][hilbertcurve]{sec}.
\end{enumerate}
Aucun de ces trois résultats ne requiert
l'axiome du choix. Nous les considérons
aujourd'hui comme de beaux résultats
mathématiques, certes surprenants.
Peut-être devrions-nous adopter la
même attitude envers le paradoxe
de Banach--Tarski.

Une précision technique s'impose
néanmoins, mais il suffit de garder
la tête froide. Les déplacements
rigides préservent le volume.
Par conséquent, les cinq\footnote{
Nous avons énoncé le paradoxe avec
«~un nombre fini de parties~».
En fait, \citet{Robinson1947} a montré
que la décomposition peut se faire
en \emph{cinq} parties, mais pas moins.
Pour une démonstration, voir
\citet[pp.~66--7]{Wagon2016}.}
parties obtenues en décomposant
la boule ne peuvent pas être
\emph{toutes mesurables}.
En gros, on ne peut donc attribuer
un volume au sens usuel à chacune
de ces parties. Il faut les imaginer
comme des «~éparpillements infinis~»
de points impossibles à représenter.
Il est peut-être «~étrange~»
de concevoir de tels ensembles
«~infiniment dispersés~».
Mais leur existence semble découler
de l'injonction, exprimée par
\stagesacc{}, de former \emph{toutes
les collections possibles}
d'ensembles déjà disponibles.

Si tout cela ne convainc pas,
voici un dernier argument,
extrinsèque, en faveur de
l'acceptation du paradoxe
de Banach--Tarski : il entraîne
immédiatement la meilleure
plaisanterie mathématique
de tous les temps.
\begin{enumerate}
	\item[] \emph{Question}.
	Quelle est l'anagramme de
	«~Banach--Tarski~» ?
	\item[] \emph{Réponse}.
	«~Banach--Tarski Banach--Tarski~».
\end{enumerate}

\end{document}
